% Cannabidiol-Alzheimer's Disease Neuroprotection — Computational Study
% PLOS Computational Biology Format
% Version v3: May 3, 2026
% Frozen baseline: model sha256 a32bf9a8b2dd2afc (git tag v3-manuscript-frozen),
%                  engine HEAD 7fd06fd8 (CPU-fork in pool, snapshot decimation).

\documentclass[10pt,letterpaper]{article}
\usepackage[top=0.85in,left=2.75in,footskip=0.75in]{geometry}

\usepackage{amsmath,amssymb,amsfonts,amsthm}
\usepackage{changepage}
\usepackage{textcomp,marvosym}
\usepackage{cite}
\usepackage{nameref,hyperref}
\usepackage[right]{lineno}
\usepackage[nopatch=eqnum]{microtype}
\DisableLigatures[f]{encoding = *, family = * }
\usepackage[table]{xcolor}
\usepackage{array}
\usepackage{verbatim}
\usepackage{booktabs}
\usepackage{graphicx}
\usepackage{subcaption}
\usepackage{tikz}
\usetikzlibrary{petri,arrows,positioning,shapes.geometric}

\newtheorem{theorem}{Theorem}
\newtheorem{lemma}[theorem]{Lemma}
\newtheorem{definition}{Definition}
\newtheorem{proposition}{Proposition}
\newtheorem{corollary}[theorem]{Corollary}
\newtheorem{example}{Example}

\newcommand{\etal}{\textit{et al.}}

\setlength{\parindent}{0pt}
\setlength{\parskip}{6pt}

\newcommand{\todo}[1]{\textcolor{red}{[TODO: #1]}}
\newcommand{\note}[1]{\textcolor{blue}{[NOTE: #1]}}
\newcommand{\pending}[1]{\textcolor{red}{[PENDING: #1]}}

\begin{document}
\vspace*{0.2in}

\begin{flushleft}
{\Large
\textbf{In Silico Mapping of Cannabidiol Neuroprotection in Alzheimer's Disease: A Stochastic Petri Net Reveals a Dissociated Therapeutic Window and a Glutathione-Mediated Bistable Cusp at Moderate Disease Severity}
}
\newline
\\
Eug\^enio Sim\~ao$^{1,*}$
\\
\bigskip
\textbf{1} Department of Computer Science, Universidade Federal de Santa Catarina (UFSC), Ararangu\'a, Santa Catarina, 88906-072, Brazil
\\
\bigskip

* eugenio.simao@ufsc.br

\end{flushleft}

% ============================================================================
% ABSTRACT
% ============================================================================
\section*{Abstract}

Multi-target pharmacology for Alzheimer's disease faces a combinatorial challenge: cannabidiol engages molecular targets that intersect the inflammatory, amyloidogenic, oxidative, and neurodegenerative cascades simultaneously, and the experimental space of dose, severity, and biomarker is too large to map purely in vivo. We present an in silico laboratory built on a hybrid stochastic--continuous biological Petri net (forty-three places, forty-eight transitions, one hundred and four arcs, and forty-three discrete events) that couples six cannabidiol targets to four disease cascades and is integrated over a clinically meaningful seven-day horizon under a twenty-seven-event maintenance dosing protocol. The model is parameterised through an explicit experiment-plan layer (parameter places for maintenance dose, loading dose, dosing interval, and disease severity) that is cleanly separated from the biological topology, so that protocols may be swept without altering the underlying network. Three findings emerge from sweeps run against a single frozen baseline. First, the untreated disease cascade is monotone in severity across every biomarker we measured: nuclear factor $\kappa$B p65 rises from $0.65$ to $2.35$ units, amyloid-$\beta$ oligomer load from $8.8$ to $117$, and reduced glutathione is depleted in proportion from $80$ to $5$ across the disease-severity grid, with microglial conservation preserved exactly. Second, cannabidiol suppresses nuclear factor $\kappa$B according to a Hill function with half-maximal inhibitory concentration of $0.08$ marking units at both mild and severe disease and Hill coefficient near $1.65$, with amyloid-$\beta$ oligomer cleared completely by the seventh day at every non-zero maintenance dose. Translated to plasma equivalents, the operative dose lies an order of magnitude below the one-to-thirty micromolar range commonly tested in preclinical studies. Third, the system exhibits an intrinsic, glutathione-mediated bistable cusp on the oxidative axis whose location depends on treatment: in the untreated cascade the bifurcation lives at moderate-to-severe disease (severity two to three, replicate coefficient of variation up to one hundred and eighteen per cent on reactive oxygen species), whereas under maintenance cannabidiol it migrates to mild disease (severity one, coefficient of variation thirty to forty-five per cent). This treatment-induced migration of the bifurcation, rather than its abolition, is the central qualitative prediction of the model: cannabidiol shifts which patients are basin-stable but cannot eliminate basin selection, motivating co-therapy targeting glutathione regeneration and stratification of clinical cohorts by oxidative-stress biomarkers in addition to amyloid load.

% ============================================================================
% AUTHOR SUMMARY
% ============================================================================
\section*{Author Summary}

Cannabidiol acts on many molecular targets at once, which makes it an attractive candidate against Alzheimer's disease but very hard to test experimentally --- there are too many doses, disease stages, and biomarkers to combine in the laboratory. We built a computational model that simulates cannabidiol treatment over a full week of realistic dosing, capturing how the drug engages inflammatory, amyloid, oxidative-stress, and neurodegeneration pathways simultaneously. The model predicts that cannabidiol suppresses brain inflammation along a sharp dose-response curve that saturates at concentrations roughly ten times below those typically tested in cell culture, and that under sustained treatment it clears amyloid-$\beta$ aggregates from the simulated brain over one week. The model further predicts that the brain's antioxidant defences sit at a tipping point whose location depends on whether the patient is being treated. In an untreated brain the tipping point lies at moderate-to-severe disease: identical patients with the same severity can fall into either a protected or a damaged oxidative state. Cannabidiol does not abolish this tipping behaviour --- it shifts the tipping point toward milder disease, so that the patients who become basin-stable under treatment are not the same patients who were stable without it. Glutathione --- the brain's principal small-molecule antioxidant --- is identified as the controlling lever. The model therefore predicts that cannabidiol monotherapy alone cannot eliminate patient-to-patient variability in oxidative outcome, that cohort studies should expect bimodal rather than unimodal oxidative-biomarker distributions in mild treated patients and in moderate untreated patients, and that pairing cannabidiol with a glutathione-restoring co-therapy should close the residual oxidative variability.

\linenumbers

% ============================================================================
% NOMENCLATURE
% ============================================================================
\section*{Nomenclature}

To preserve readability we use the full biological name on first mention and adopt the abbreviation sparingly thereafter. Table~\ref{tab:nomenclature} maps each abbreviation used in this paper to its corresponding biological compound and pathway role.

\begin{table}[htbp]
\centering
\small
\caption{Abbreviations used in this work.}
\label{tab:nomenclature}
\begin{tabular}{l p{0.55\textwidth} l}
\toprule
Abbreviation & Compound / pathway element & Role \\
\midrule
CBD & Cannabidiol & Multi-target phytocannabinoid (the drug) \\
AD  & Alzheimer's disease & Pathological context \\
A$\beta$ & Amyloid-$\beta$ peptide (monomer, oligomer, plaque) & Pathological aggregate \\
APP & Amyloid precursor protein & Substrate for $\gamma$-secretase \\
NF$\kappa$B & Nuclear factor $\kappa$-light-chain-enhancer of activated B cells (subunit p65) & Master inflammatory transcription factor \\
I$\kappa$B  & Inhibitor of NF$\kappa$B & Cytoplasmic NF$\kappa$B chaperone \\
IKK         & I$\kappa$B kinase & Activator of NF$\kappa$B release \\
TNF$\alpha$, IL-1$\beta$, IL-6 & Tumour necrosis factor $\alpha$, interleukin-1$\beta$, interleukin-6 & Pro-inflammatory cytokines \\
COX-2 & Cyclooxygenase-2 & Inflammatory enzyme, NF$\kappa$B target \\
Nrf2 & Nuclear factor erythroid 2-related factor 2 & Master antioxidant transcription factor \\
Keap1 & Kelch-like ECH-associated protein 1 & Cytoplasmic Nrf2 anchor \\
HO-1  & Heme oxygenase 1 & Antioxidant enzyme, Nrf2 target \\
SOD   & Superoxide dismutase & Antioxidant enzyme \\
ROS   & Reactive oxygen species & Oxidative driver \\
GSH, GSSG & Reduced and oxidised glutathione & Antioxidant redox couple \\
BDNF  & Brain-derived neurotrophic factor & Neurotrophic survival signal \\
GPR3  & G-protein-coupled receptor 3 & Modulator of $\gamma$-secretase activity \\
PPAR$\gamma$ & Peroxisome proliferator-activated receptor $\gamma$ & Nuclear anti-inflammatory receptor \\
5-HT$_{1A}$  & Serotonin 1A receptor & Neurotrophic G-protein-coupled receptor \\
A$_{2A}$    & Adenosine A$_{2A}$ receptor & Microglial polarisation receptor \\
NAC & N-acetylcysteine & Glutathione precursor (proposed co-therapy) \\
\bottomrule
\end{tabular}
\end{table}

% ============================================================================
% INTRODUCTION
% ============================================================================
\section{Introduction}
\label{sec:introduction}

Alzheimer's disease accounts for the majority of dementia cases worldwide, affecting more than fifty-five million people and projected to surpass one hundred and thirty-nine million by 2050~\cite{who2023}. Disease-modifying therapies remain elusive. Amyloid-targeting monoclonal antibodies such as lecanemab and donanemab show only modest clinical benefit and carry significant adverse effects~\cite{vandyck2023}, and over twenty-five trials of anti-inflammatory monotherapies have consistently failed to slow cognitive decline despite reducing inflammatory biomarkers~\cite{adapt2007,imbimbo2010}. These repeated negative results highlight two deeper problems: the pathology is driven by interacting cascades that no single-target drug can fully suppress, and the experimental design space --- spanning dose, dosing schedule, disease severity, patient subtype, and biomarker --- is too large to map purely in vivo.

Cannabidiol, a non-psychoactive phytocannabinoid, has emerged as a candidate multi-target agent with documented activity at six molecular targets relevant to this disease. Preclinical evidence demonstrates inverse agonism of G-protein-coupled receptor 3, with consequent reduction of $\gamma$-secretase activity and amyloid-$\beta$ production~\cite{huang2015}; nuclear activation of peroxisome proliferator-activated receptor $\gamma$, with consequent suppression of inflammation driven by nuclear factor $\kappa$B~\cite{esposito2011}; dissociation of the cytoplasmic Keap1--Nrf2 complex, releasing nuclear factor erythroid 2-related factor 2 to drive transcription of antioxidant response elements~\cite{casares2020}; agonism of the serotonin 1A receptor, promoting neuroprotection through brain-derived neurotrophic factor~\cite{russo2005}; agonism of the adenosine A$_{2A}$ receptor, driving anti-inflammatory polarisation of microglia~\cite{mecha2013}; and direct modulation of mitochondrial reactive oxygen species through the Keap1--Nrf2 axis~\cite{booz2011}. No published computational framework, however, integrates these six targets with the four pathological cascades over a clinically realistic time horizon to predict therapeutic windows, dose--response curves, and the emergent multi-stable phenomena that arise from cross-cascade coupling.

The fundamental question posed by this multiplicity is not whether to experiment, but which experiments to prioritise --- which dose ranges contain critical transitions, which biomarkers must be measured independently, and which patient strata require adjusted dosing or co-therapy. A computational model can serve as an in silico laboratory in this respect, mapping the therapeutic landscape ahead of wet-lab commitment, identifying the most informative experimental conditions, and specifying the dose ranges and biomarker panels that will yield the sharpest discrimination between candidate outcomes. Existing models, however, address components of the cannabidiol--Alzheimer system in isolation. Amyloid kinetic models simulate aggregation without pharmacological intervention~\cite{proctor2010,kyrtsos2015}; nuclear factor $\kappa$B signalling models capture oscillatory dynamics without amyloid or cannabinoid input~\cite{hoffmann2002}; Nrf2 pathway models characterise the antioxidant response outside any disease context~\cite{zhang2017}; cannabidiol pharmacokinetic models describe absorption and metabolism without mechanistic targets~\cite{zgair2016,millar2018}; and quantitative systems-pharmacology platforms for Alzheimer's address cholinergic and glutamatergic systems but not the targets engaged by cannabidiol~\cite{geerts2016}. This fragmentation prevents prediction of the very phenomena that confound clinical practice, such as the dissociation between inflammatory efficacy and neuroprotection that has plagued anti-inflammatory trials~\cite{heneka2015}, or the bistable redox dynamics that may underlie patient-to-patient variability~\cite{cassano2020,watt2017}.

This work addresses the integration gap with a clinically realistic model that strictly separates biology from protocol. The biological topology --- forty-three places, forty-eight transitions, and one hundred and four arcs --- couples cannabidiol's six molecular targets to the four cascades; an experiment-plan layer of parameter places (maintenance dose, loading dose, dosing interval, disease severity) and forty-three discrete events delivers the dosing protocol and installs disease pathology without modifying the underlying network. Sweeping this layer over a seven-day horizon recovers three findings reported in this paper. The first is the verification step: in the absence of cannabidiol, increasing disease severity drives a strict-monotone cascade response across every measured biomarker (inflammation up, amyloid load up, glutathione down, microglia activated), with the topological conservation invariants preserved exactly --- confirming that the experiment-plan layer perturbs only the biological initial condition and that subsequent findings are emergent properties of the topology. The second is therapeutic: cannabidiol suppresses nuclear factor $\kappa$B p65 along a Hill curve with sub-micromolar half-maximal inhibitory concentration weakly modulated by disease severity, and clears amyloid-$\beta$ oligomer to zero over the seven-day horizon at every non-zero maintenance dose. The third is qualitative: the system carries an intrinsic glutathione-mediated bistable cusp on the oxidative axis, located at moderate-to-severe disease in the untreated cascade and at mild disease under treatment. Cannabidiol therefore does not abolish basin selection on the redox axis but shifts which patients sit on the cusp --- a treatment-induced migration of the bifurcation that motivates the central translational prediction of the work, that effective therapy at any severity stratum requires a co-therapy targeting glutathione regeneration, with N-acetylcysteine the natural candidate, and that future trials should stratify by oxidative-stress biomarkers in addition to amyloid load.

% ============================================================================
% METHODS
% ============================================================================
\section{Methods}
\label{sec:methods}

\subsection{Computational Framework}

The model is implemented within the Signal Hierarchical Petri Net framework~\cite{simao2026shypn}, a hybrid stochastic--continuous simulator for biological Petri-net models. The engine couples tau-leaping stochastic simulation~\cite{gillespie2001tauleap} with deterministic ordinary-differential-equation integration via Runge--Kutta in an operator-split scheme: per simulation step, continuous flow is integrated, then stochastic firings are sampled, then the marking is finalised. The arc semantics distinguish four roles --- normal arcs that transfer mass, test arcs that perform non-consuming presence checks, inhibitor arcs whose presence disables a transition, and signal-flow arcs that carry mass under cascade gating with a basin-floor enabling rule --- which together implement the signal hierarchy that ranks transitions by control depth.

The biological topology and the experimental protocol are kept architecturally separate. The object-net carries the biology and its dynamics emerge entirely from its own topology; the experiment-plan layer comprises four parameter places (maintenance dose, loading dose, dosing interval, and disease severity) and a set of discrete events. Parameter places do not appear in any biological rate function and have no arcs to or from biological places; the only legal bridge between layers is event-mediated, in which a discrete event reads parameter values and applies a marking change to a biological place at scheduled times. This discipline makes protocol sweeps non-invasive: the same biological topology is reused for every condition, and the conservation invariants on the biological places (described below) cannot be broken by editing the experimental design.

All numerical results in this paper were produced from the model file \texttt{canabidiol-q1-testable.shy} with sha256 prefix \texttt{a32bf9a8b2dd2afc} (git tag \texttt{v3-manuscript-frozen}) and engine commit \texttt{7fd06fd8}. Each sweep run directory carries a \texttt{provenance.json} (recording client and server git commit, branch, dirty-tree status, and dispatch timestamp) and a byte-identical \texttt{model\_snapshot.shy} (the exact bytes the worker loaded), making every reported number independently reconstructible regardless of subsequent edits to the repository. Sweeps are dispatched through a pool of CPU-fork workers with snapshot auto-decimation enabled.

\subsection{Model Architecture}

The model comprises forty-three places, forty-eight transitions, one hundred and four arcs, and forty-three discrete events. Of the events, fifteen install disease pathology immediately after $t = 0$, one delivers the loading dose at $t = 0$, and twenty-seven deliver maintenance doses at six-hour cadence over the seven-day horizon. Of the places, thirty-nine are biological and four constitute the experiment-plan layer.

The biological places fall into seven functional groups. Pharmacokinetics is represented by extracellular cannabidiol, plasma cannabidiol, intracellular cannabidiol, and a clearance compartment. The cannabidiol target group contains active and inactive forms of G-protein-coupled receptor 3, the active forms of peroxisome proliferator-activated receptor $\gamma$, the serotonin 1A receptor, and the adenosine A$_{2A}$ receptor. The amyloid cascade contains amyloid precursor protein messenger RNA, amyloid precursor protein, $\gamma$-secretase, and three amyloid-$\beta$ aggregation states (monomer, oligomer, plaque). The neuroinflammation group contains the cytoplasmic NF$\kappa$B--I$\kappa$B complex and free NF$\kappa$B p65, the I$\kappa$B kinase, the cytokines tumour necrosis factor $\alpha$, interleukin-1$\beta$, and interleukin-6, the inflammatory enzyme cyclooxygenase-2, and the M1 and M2 microglial polarisation states. The oxidative-stress group contains the cytoplasmic Keap1--Nrf2 complex, free Nrf2, heme oxygenase 1, superoxide dismutase, reactive oxygen species, and the reduced/oxidised glutathione couple. The neurodegeneration group contains a single neuron-health state and brain-derived neurotrophic factor. Finally, four environmental compartment places (temperature, pH, age, and an age-derived factor) encode background physiological context.

The biological topology enforces four place conservation invariants:
\begin{align}
\text{NF}\kappa\text{B--I}\kappa\text{B} + \text{NF}\kappa\text{B p65} &= \text{const}, &
\text{M1} + \text{M2} &= \text{const}, \\
\text{Keap1--Nrf2} + \text{free Nrf2} &= \text{const}, &
\text{GSH} + \text{GSSG} &= \text{const}.
\end{align}
These invariants constrain the system's phase space so that mass is redistributed between active and inactive forms rather than created or destroyed. The glutathione invariant, in particular, is critical for the bistability result of Section~\ref{sec:cusp}: depletion of reduced glutathione under disease load is matched by accumulation of oxidised glutathione, and the antioxidant capacity coefficient depends linearly on the reduced fraction.

Disease pathology is installed by fifteen events of the schematic form ``$\text{species} := \text{species} + \text{disease severity} \cdot \delta$'' that fire once at $t \approx 0$. They install severity-proportional amyloid-$\beta$ at all three aggregation stages; elevation of the cytokine pool (tumour necrosis factor $\alpha$, interleukin-1$\beta$, interleukin-6, cyclooxygenase-2); upregulation of I$\kappa$B kinase and free NF$\kappa$B p65; M1 polarisation of microglia; elevation of reactive oxygen species and depletion of reduced glutathione; and upregulation of amyloid precursor protein and $\gamma$-secretase. A disease-severity value of zero leaves the model in its healthy initial marking; values of one, two, and five correspond, by construction, to mild, moderate, and severe stages.

The dosing protocol has a similarly compact specification: a single loading-dose event at $t = 0$ adds the loading-dose parameter to extracellular cannabidiol, and twenty-seven maintenance-dose events fire at six-hour intervals across the seven-day horizon, each adding the maintenance-dose parameter to the same place. The protocol approximates four-times-daily oral cannabidiol in a chronic regimen. A three-compartment pharmacokinetic submodel (plasma $\leftrightarrow$ extracellular $\leftrightarrow$ intracellular, with linear clearance) links the loaded extracellular pool to the intracellular pool that the intracellular cannabidiol targets read. Q$_{10}$ temperature dependence and a pH-acidosis modifier are applied to enzymatic transitions to keep environmental context coupled to enzyme kinetics.

The six cannabidiol mechanisms correspond to six dedicated transitions: inverse agonism of G-protein-coupled receptor 3 (depleting the active form, opposed by recycling); intracellular activation of peroxisome proliferator-activated receptor $\gamma$, which then suppresses NF$\kappa$B p65; dissociation of the Keap1--Nrf2 complex, dual-driven by reactive oxygen species and intracellular cannabidiol; serotonin 1A-mediated production of brain-derived neurotrophic factor; adenosine A$_{2A}$-mediated M1$\to$M2 polarisation; and direct reactive-oxygen-species modulation embedded in the Keap1--Nrf2 dissociation rate. These six transitions are the points at which the experiment-plan layer (via dosed extracellular cannabidiol, redistributed by the pharmacokinetic submodel) reaches into the biology.

The single transition labelled ``Neurotoxicity'' integrates oligomer-, oxidative-, and cytokine-driven neuronal death. Its rate function reads
\begin{align}
r_{\text{Neurotox}} = \big( & 0.003 \cdot \tfrac{\text{A}\beta_{\text{olig}}}{10 + \text{A}\beta_{\text{olig}}}
+ 0.004 \cdot \tfrac{\max(0, \text{ROS}-1)}{15 + \max(0, \text{ROS}-1)} \nonumber \\
+ & 0.003 \cdot \tfrac{\max(0, \text{TNF}\alpha-0.5)}{10 + \max(0, \text{TNF}\alpha-0.5)} \big) \cdot Q_{10} \cdot (1 + 0.4 \cdot \text{pH}_{\text{acid}}) \cdot \text{Age}_{\text{f}} .
\end{align}
The Hill thresholds (one for reactive oxygen species, one half for tumour necrosis factor $\alpha$) implement basal-tolerance behaviour: small homeostatic levels do not drive damage. The ROS coefficient of $0.004$ is the value used throughout this manuscript and in all reported sweeps. An alternative softening (a coefficient of $0.001$ with threshold raised to twenty) was investigated in pilot work and shown to flip the model's global attractor; this observation provides the empirical demonstration of multi-stability discussed in Section~\ref{sec:disc-bistability} and motivates the discipline of freezing the model bytes for the entire manuscript.

\subsection{Experimental Design}

Three sweeps were dispatched against the frozen baseline. Each sweep varies parameter places only, leaves the biological topology untouched, and runs thirty stochastic replicates per condition with unique seeds, a seven-day horizon ($t = 604\,800$ s), tau-leaping with $\tau$-epsilon of $0.03$, and maximum tau of $0.1$.

The therapeutic-window factorial sweep, labelled ``Q4r-final'' in the run archive, varies maintenance dose over $\{0, 0.5, 2, 5\}$ and disease severity over $\{0, 1, 2, 5\}$, giving sixteen factorial cells plus an explicit baseline cell at the model defaults. Loading dose is fixed at ten units and dosing interval at six hours. The total trajectory budget is five hundred and ten replicates.

The inflammatory dose-response refinement, labelled ``Q5-final'', varies maintenance dose over a finer grid $\{0, 0.05, 0.1, 0.2, 0.35, 0.5\}$ at two disease severities (one and five), plus a baseline. The intermediate severity of two is omitted because the bistable cusp identified in the Q4r-final sweep would require at least forty replicates per cell to resolve cleanly and is not required for the half-maximal inhibitory concentration measurement. The total budget is three hundred and ninety replicates.

The disease cascade monotonicity sweep, labelled ``Q3-final'', sweeps disease severity over $\{0, 0.5, 1, 2, 3, 5\}$ in the absence of cannabidiol, plus a baseline. It verifies that the disease-installation events produce a strict-monotone cascade response across severity in the untreated topology and that the conservation invariants are preserved across all severity levels. The total budget is two hundred and ten replicates.

Endpoint statistics (mean, standard deviation, minimum, and maximum across each replicate cohort), inter-replicate variance, bimodality detection on key endpoints, and per-condition Hill fits of the form $r = r_{\max} \cdot \text{IC}_{50}^n / (\text{IC}_{50}^n + [\text{CBD}]^n)$ on the inflammatory readouts are computed from the per-condition CSV output of the sweep runner.

% ============================================================================
% RESULTS
% ============================================================================
\section{Results}
\label{sec:results}

All numerical values reported below derive from the three frozen-baseline sweeps Q4r-final (run \texttt{20260503\_130113}, sixteen factorial cells plus baseline), Q5-final (run \texttt{20260503\_133835}, twelve factorial cells plus baseline), and Q3-final (run \texttt{20260503\_140103}, six severity cells plus baseline), each with thirty replicates per cell on the model bytes archived under git tag \texttt{v3-manuscript-frozen}. Per-condition CSV summaries and per-replicate trajectories are archived under \texttt{workspace/projects/canabidiol/experiments/results/} alongside the run-local \texttt{provenance.json} and \texttt{model\_snapshot.shy} for each sweep.

\subsection{Disease Cascade Monotonicity Across Severity}
\label{sec:r-q3}

We first verify, in the absence of cannabidiol, that the experiment-plan layer drives the expected biological cascade across the severity axis. Increasing disease severity from zero to five produces a monotone increase in nuclear factor $\kappa$B p65 from $0.65$ to $2.35$ units, a $13.4$-fold monotone increase in amyloid-$\beta$ oligomer load from $8.8$ to $117.4$, a $5.5$-fold rise in plaque load from $0.64$ to $3.49$, and a $16$-fold monotone depletion of reduced glutathione from $80$ to $5$ units. M1 microglial activation rises from $2$ to $7$ tokens while the M2 pool falls from $43$ to $63$ in opposite proportion, leaving the microglial conservation invariant preserved exactly at every severity level (sums of forty-three to seventy tokens that depend on the severity-installed quantum but show zero replicate variance once installed). The cascade values are summarised in Table~\ref{tab:q3-cascade}. The reactive-oxygen-species axis is non-monotone in mean by design --- the antioxidant capacity coefficient depends on the depleted glutathione pool through an enzymatic floor, and the resulting bistability is taken up in Section~\ref{sec:cusp}. The deterministic glutathione trajectory, the exact microglial conservation, and the order-of-magnitude amyloid-oligomer dynamic range across the severity grid jointly confirm that the disease-installation events perturb only the biological initial condition and that the qualitative responses we attribute to cannabidiol in the following subsections are emergent properties of the topology itself.

\begin{table}[htbp]
\centering
\caption{Q3-final cascade response across disease severity in the untreated cascade. Mean endpoints over thirty replicates; standard deviation in parentheses; \texttt{run\_20260503\_140103}.}
\label{tab:q3-cascade}
\begin{tabular}{lrrrrrr}
\toprule
Disease severity & 0 & 0.5 & 1 & 2 & 3 & 5 \\
\midrule
NF$\kappa$B p65         & 0.65 (0.01) & 0.81 (0.01) & 1.14 (0.03) & 1.64 (0.11) & 1.98 (0.13) & 2.35 (0.02) \\
A$\beta$ oligomer       & 8.8 (0.5)   & 13.7 (0.8)  & 26.1 (1.1)  & 50.8 (2.5)  & 90.3 (4.2)  & 117.4 (5.8) \\
A$\beta$ plaque         & 0.64 (0.55) & 0.78 (0.66) & 1.42 (0.74) & 2.64 (1.12) & 3.38 (1.44) & 3.49 (1.10) \\
Reduced glutathione     & 80.0        & 72.5        & 65.0        & 50.0        & 35.0        & 5.0         \\
Reactive oxygen species & 0           & 0           & 65.9 (0.3)  & 26.9 (31.7) & 50.9 (36.7) & 54.9 (0.1)  \\
Microglia M1            & 2.0 (1.3)   & 3.0 (2.1)   & 4.4 (2.2)   & 4.4 (2.3)   & 6.5 (3.4)   & 7.0 (2.6)   \\
Microglia M2            & 43.0 (1.3)  & 44.5 (2.1)  & 45.6 (2.2)  & 50.6 (2.3)  & 53.5 (3.4)  & 63.2 (2.6)  \\
\bottomrule
\end{tabular}
\end{table}

\subsection{The Therapeutic Surface Across Dose and Severity}
\label{sec:r-q4r}

The Q4r-final sweep maps the four-by-four therapeutic surface over maintenance dose and disease severity. Three readouts characterise the response: nuclear factor $\kappa$B p65 captures the inflammatory axis, amyloid-$\beta$ oligomer load captures the amyloidogenic axis, and reactive oxygen species captures the redox axis (Figure~\ref{fig:q4r-surface}, Table~\ref{tab:q4r-surface}).

The inflammatory surface is steep and saturating. In the absence of treatment, NF$\kappa$B p65 rises monotonically with severity from $1.14$ at severity zero to $2.17$ at severity five; once any non-zero maintenance dose is applied the entire severity axis is suppressed by approximately a factor of three or more, settling between $0.30$ and $0.66$ across the cell grid, with no further dependence on dose between maintenance levels of $0.5$ and $5$. The Q4r-final factorial therefore brackets the inflammatory inhibitory concentration in the open interval $(0, 0.5)$ and motivates the refined Q5-final sweep (Section~\ref{sec:r-q5}).

The amyloidogenic axis exhibits the strongest treatment effect of any biomarker in the sweep: amyloid-$\beta$ oligomer load is cleared to zero at the seven-day endpoint in every cell of the factorial, including the no-treatment column where the same biomarker plateaued at $14$ to $117$ units in the untreated Q3-final cascade (Section~\ref{sec:r-q3}). The clearance is driven by the inflammation-mediated downregulation of amyloid precursor protein once peroxisome proliferator-activated receptor $\gamma$ activation begins to suppress NF$\kappa$B p65; even the residual NF$\kappa$B activity in the no-treatment column of the Q4r-final factorial is sufficient, over the seven-day horizon, to halt new oligomer formation while existing oligomers progress to plaque or are degraded. Note that this treatment-driven oligomer clearance is a property of the seven-day horizon: it does not appear in the shorter pilot integrations of earlier work and is one of the cleanest predictions to test against in vivo amyloid-imaging time courses.

The redox axis is qualitatively different: it is dose-insensitive but severity-shaped, and at one severity exhibits a bistable cusp that is taken up in Section~\ref{sec:cusp}. The redox response shows essentially no dose dependence within each severity row of the factorial (the four maintenance doses of $0$, $0.5$, $2$, $5$ produce statistically indistinguishable mean reactive-oxygen-species endpoints), confirming that cannabidiol does not act directly on the redox axis in this topology --- its only redox-relevant arc is the inflammation-mediated cytokine route, which is already saturated by the lowest non-zero dose tested.

\begin{table}[htbp]
\centering
\caption{Q4r-final endpoint surface across maintenance dose and disease severity. Mean over thirty replicates; standard deviation in parentheses; \texttt{run\_20260503\_130113}. Amyloid-$\beta$ oligomer load is identically zero across the entire factorial at $t = 7$~d.}
\label{tab:q4r-surface}
\begin{tabular}{lrrrrr}
\toprule
 & \multicolumn{4}{c}{Maintenance dose} & \\
\cmidrule(lr){2-5}
Severity & 0 & 0.5 & 2 & 5 & Readout \\
\midrule
0 & 1.14 (0.001) & 0.30 (0.000) & 0.30 (0.000) & 0.30 (0.000) & NF$\kappa$B p65 \\
1 & 1.35 (0.024) & 0.46 (0.034) & 0.45 (0.051) & 0.45 (0.041) & NF$\kappa$B p65 \\
2 & 1.59 (0.005) & 0.52 (0.003) & 0.52 (0.002) & 0.52 (0.003) & NF$\kappa$B p65 \\
5 & 2.17 (0.000) & 0.66 (0.000) & 0.66 (0.000) & 0.66 (0.000) & NF$\kappa$B p65 \\
\midrule
0 & 0    (0)    & 0    (0)    & 0    (0)    & 0    (0)    & Reactive O species \\
1 & 51.1 (17)   & 53.0 (14)   & 47.4 (22)   & 51.1 (17)   & Reactive O species \\
2 & 72.0 (1.3)  & 71.6 (0.91) & 71.6 (0.77) & 71.7 (0.83) & Reactive O species \\
5 & 46.7 (0.002)& 46.7 (0)    & 46.7 (0)    & 46.7 (0)    & Reactive O species \\
\bottomrule
\end{tabular}
\end{table}

\begin{figure}[htbp]
\centering
\includegraphics[width=\textwidth,draft]{fig_q4r_surface_placeholder.pdf}
\caption{\pending{Q4r-final therapeutic surface.} Three heatmaps over maintenance dose and disease severity: (a) NF$\kappa$B p65 (saturating inhibition for any non-zero maintenance dose), (b) reactive oxygen species (severity-shaped, dose-insensitive, bistable at severity one), (c) amyloid-$\beta$ oligomer (uniformly zero at $t = 7$~d across the entire factorial). The redox bistability at severity one and its migration relative to the untreated cascade are the central qualitative finding of this paper (Section~\ref{sec:cusp}).}
\label{fig:q4r-surface}
\end{figure}

The pharmacokinetic submodel responds linearly to maintenance dose, as shown in Table~\ref{tab:q4r-pk}: both the plasma and intracellular cannabidiol endpoints scale exactly proportionally over the four maintenance levels (plasma cannabidiol per maintenance unit is $0.376$, intracellular cannabidiol per maintenance unit is $0.687$, both constant to three decimal places across the factorial). This linear PK response confirms that observed endpoint differences across cells reflect a pharmacologically meaningful dose gradient rather than a saturated or starved compartment.

\begin{table}[htbp]
\centering
\caption{Q4r-final endpoint pharmacokinetics at disease severity one. Mean values at $t = 7$~d over thirty replicates.}
\label{tab:q4r-pk}
\begin{tabular}{rrr}
\toprule
Maintenance dose & Plasma cannabidiol & Intracellular cannabidiol \\
\midrule
0    & 0.000 & 0.000 \\
0.5  & 0.188 & 0.343 \\
2    & 0.752 & 1.373 \\
5    & 1.880 & 3.433 \\
\bottomrule
\end{tabular}
\end{table}

\subsection{A Sub-Micromolar Half-Maximal Inhibitory Concentration for Inflammation}
\label{sec:r-q5}

The Q4r-final factorial saturates the inflammatory readout already at a maintenance dose of $0.5$ units (Table~\ref{tab:q4r-surface}), placing the half-maximal inhibitory concentration in the open interval $(0, 0.5)$. The Q5-final refinement at six maintenance levels in this interval yields the dose--response curves of Figure~\ref{fig:q5-hill}. A four-parameter Hill fit gives a half-maximal inhibitory concentration of $0.080$ marking units with Hill coefficient $n = 1.64$ at mild severity (sum of squared residuals $0.078$), and $0.077$ marking units with $n = 1.65$ at severe disease (sum of squared residuals $0.134$). The two fits are essentially superimposed on one another, with the mild- and severe-severity curves differing principally in their unsuppressed plateau ($r_0 = 3.17$ vs.\ $4.43$ NF$\kappa$B-p65 units) rather than in the position or steepness of the half-saturation: the inflammatory dose-response is severity-independent in this regime. In plasma-equivalent units this places the operative dose in the sub-micromolar range, well below the doses commonly tested in preclinical cannabidiol studies, which typically span one to thirty micromolar. The implication --- discussed below --- is that the historical preclinical literature has likely operated at or near saturation of the very inflammatory pathway it has been measuring.

\begin{figure}[htbp]
\centering
\includegraphics[width=0.7\textwidth,draft]{fig_q5_hill_placeholder.pdf}
\caption{\pending{Q5-final Hill fit.} NF$\kappa$B p65 endpoint (mean and standard deviation across thirty replicates) versus maintenance dose on the refined grid $\{0, 0.05, 0.1, 0.2, 0.35, 0.5\}$ at two disease severities. Curves are four-parameter Hill fits; the half-maximal inhibitory concentration ($\text{IC}_{50} = 0.080$ at mild, $0.077$ at severe disease) and Hill coefficient ($n \approx 1.65$ at both severities) are annotated.}
\label{fig:q5-hill}
\end{figure}

\subsection{A Glutathione-Mediated Bistable Cusp Whose Location Depends on Treatment}
\label{sec:cusp}

The redox axis exhibits an intrinsic bistable cusp whose location is itself a function of treatment. In the untreated cascade (Q3-final, Table~\ref{tab:q3-cascade}, no cannabidiol), the bifurcation lies at moderate-to-severe disease: at severity two the per-replicate reactive-oxygen-species endpoint distribution is broad and bimodal (mean $26.9$, standard deviation $31.7$, coefficient of variation $118\%$), at severity three it is similarly bimodal but biased toward the high basin (mean $50.9$, standard deviation $36.7$, coefficient of variation $72\%$), and at the surrounding severities the same statistic collapses to under one per cent. Identical conditions partition into two outcome classes, one with low reactive oxygen species and a residual reduced-glutathione pool, the other with high reactive oxygen species and depleted reduced glutathione (Figure~\ref{fig:cusp}, panel a). The cusp is therefore not a treatment artefact: it is intrinsic to the disease cascade itself.

Under maintenance cannabidiol (Q4r-final, Table~\ref{tab:q4r-surface}), the same bistable signature reappears but at a different severity. The replicate distribution is now tight at severity zero, severity two ($71.6 \pm 0.8$, coefficient of variation $1\%$), and severity five ($46.7 \pm 0$, coefficient of variation under $0.005\%$), and bimodal at severity one (means $47$ to $53$, standard deviations $14$ to $22$, coefficient of variation $30$ to $45\%$ across the four maintenance doses tested). Treatment therefore migrates the bifurcation by approximately one severity step toward milder disease without abolishing it: which patient sits on the cusp depends on whether the patient is treated, but the cusp itself remains. Because the bistable signature persists at all four maintenance doses tested at severity one with no detectable dose-modulation of its location or width (Table~\ref{tab:q4r-surface}), the migration is determined by the choice to treat rather than by the choice of dose --- once any non-zero dose suppresses the inflammatory cascade, the redox cusp settles at its treated location.

\begin{figure}[htbp]
\centering
\includegraphics[width=0.85\textwidth,draft]{fig_cusp_bimodality_placeholder.pdf}
\caption{\pending{Treatment-induced migration of the bistable cusp.} (a) Per-replicate reactive-oxygen-species endpoints across the disease-severity axis in the untreated cascade (Q3-final, no cannabidiol): tight at severities zero, one, and five; bimodal at severities two ($\text{CV} = 118\%$) and three ($\text{CV} = 72\%$). (b) The same panel under maintenance cannabidiol (Q4r-final at maintenance dose two): tight at severities zero, two, and five; bimodal at severity one ($\text{CV} \approx 45\%$), with the location of the cusp migrated approximately one severity step relative to (a). (c) Mechanistic phase portrait: antioxidant capacity $K(\text{GSH}) = a + b \cdot \text{GSH}$ as a function of reduced glutathione, intersecting the amyloid-driven reactive-oxygen-species production curve in a narrow severity window.}
\label{fig:cusp}
\end{figure}

The mechanistic origin is transparent. Two transitions govern reactive-oxygen-species dynamics: a basal production driven by oligomer load and inflammation, and an antioxidant scavenging driven by superoxide dismutase, heme oxygenase 1, and reduced glutathione. The scavenging coefficient takes the form $K(\text{GSH}) = a + b \cdot \text{GSH}$ with $a$ a constant enzymatic floor and $b$ the per-unit glutathione contribution. Disease severity depletes reduced glutathione approximately linearly through the GSH--GSSG conservation cycle ($72.5$, $65$, $50$, $35$, $5$ at the severities of Table~\ref{tab:q3-cascade}). At low severity $K$ exceeds the production driver and reactive oxygen species settle at a low fixed point; at high severity $K$ falls below the driver and reactive oxygen species settle at a high fixed point clamped by the self-inhibition term in the production rate. At one intermediate severity $K$ crosses the production curve in a narrow window where stochastic pulses can switch a trajectory between the two basins --- a textbook bistable cusp. Cannabidiol acts on reactive oxygen species only through the inflammation-mediated cytokine axis (peroxisome proliferator-activated receptor $\gamma$ to NF$\kappa$B to tumour necrosis factor $\alpha$ to oligomer-driven reactive oxygen species); by removing the inflammatory contribution to the production curve it lowers the curve uniformly across severity and shifts the crossover with $K(\text{GSH})$ to a milder severity, exactly as observed. The model therefore predicts that cannabidiol cannot eliminate basin selection on the redox axis but reassigns which severity stratum is intrinsically bimodal. The clinical implication --- that the patients who become high-variance under cannabidiol monotherapy are not the same patients who were high-variance without it, and that a glutathione-restoring co-therapy is required at any severity to abolish basin selection altogether --- is taken up in the Discussion.

% ============================================================================
% DISCUSSION
% ============================================================================
\section{Discussion}
\label{sec:discussion}

This in silico laboratory maps the therapeutic landscape of cannabidiol over a clinically realistic seven-day dosing horizon and yields three findings that together rationalise the historical pattern of clinical-trial outcomes for inflammation-axis drugs in Alzheimer's disease, identify a previously unappreciated patient-stratification axis (oxidative-stress basin), and predict a specific co-therapy strategy.

The first finding is the verification step that anchors the rest. The Q3-final cascade, run with no cannabidiol, recovers a strict-monotone severity response across every measured biomarker (NF$\kappa$B p65 from $0.65$ to $2.35$, amyloid-$\beta$ oligomer from $8.8$ to $117$, reduced glutathione depleted from $80$ to $5$, microglial pool conserved exactly) and confirms that the experiment-plan layer perturbs only the biological initial condition. The second finding is therapeutic and quantitative. Cannabidiol suppresses NF$\kappa$B p65 along a Hill curve with half-maximal inhibitory concentration of $0.080$ marking units at mild severity and $0.077$ at severe disease, with Hill coefficient near $1.65$ at both severities --- essentially severity-independent. Over the seven-day horizon the same treatment clears amyloid-$\beta$ oligomer to zero across the entire $4 \times 4$ factorial of dose and severity, including in the no-treatment column where the same biomarker plateaued at $14$ to $117$ units in the untreated cascade. The half-maximal inhibitory concentration of $0.08$ marking units is at least an order of magnitude below the doses commonly tested in cannabidiol pharmacology studies, which typically span one to thirty micromolar. If the inflammatory readout saturates by half a micromolar, then differences observed in vitro across the one-to-thirty micromolar range are likely reflecting non-inflammatory mechanisms --- off-target binding, generic antioxidant action, or cytotoxicity at the high end --- rather than the inflammatory pathway proper. Historical dose-finding studies focused on inflammatory endpoints may therefore have systematically operated above saturation of the very pathway they were measuring, and future preclinical work should use sub-micromolar dose grids to distinguish cannabidiol-specific effects from generic anti-inflammatory action.

\subsection{Bistable Redox Cusp and Patient Stratification}
\label{sec:disc-cusp}

The bistable cusp identified in Section~\ref{sec:cusp} is, to our knowledge, a novel computational prediction in the cannabidiol--Alzheimer literature, both in its existence and in its observed migration under treatment. It frames patient-to-patient variability not as pharmacokinetic noise but as basin selection in a multi-stable redox landscape, and the migration of the bifurcation by approximately one severity step toward milder disease under maintenance cannabidiol means that the patients who become high-variance under treatment are not the same patients who were high-variance without it. Two specific testable predictions follow. First, cohort studies of plasma 8-hydroxy-2'-deoxyguanosine, isoprostanes, and total glutathione should observe wider dispersion --- bimodal rather than unimodal --- in moderate-to-severe untreated patients, but in mild patients on chronic low-dose cannabidiol, with the location of the bimodal stratum shifting between the two arms. Second, co-therapy with N-acetylcysteine or another glutathione precursor should abolish basin selection altogether by lifting the antioxidant capacity coefficient $K(\text{GSH})$ above the production curve at every severity. The model predicts that N-acetylcysteine monotherapy will be insufficient on its own --- it does not address the inflammatory axis, and would leave the production curve where the untreated cascade puts it --- but that combined with cannabidiol it should produce a unimodal redox response across the entire severity grid. A two-by-two trial design (cannabidiol with or without N-acetylcysteine) at fixed mild and at fixed moderate severity should reveal a synergistic, rather than merely additive, collapse of replicate-level redox variance.

\subsection{Multi-Stability in the Computational Model: A Methodological Note}
\label{sec:disc-bistability}

In pilot work supporting this study, a single-parameter softening of the reactive-oxygen-species coefficient in the neurotoxicity rate function (changing the prefactor from $0.004$ to $0.001$ on the reactive-oxygen-species term) was observed to flip the model's global attractor across all thirty replicates, from the amyloid-disease basin (replete reduced glutathione, oligomer-positive) to a redox-collapse basin (reduced glutathione $\to 0$, oligomer cleared, neuron-health pool driven to its lower clamp). This was a methodologically important observation: the model has at least two reachable disease attractors at the same parameter values, and which one is sampled depends on early-trajectory reactive-oxygen-species dynamics, themselves sensitive to the rate-coefficient structure. We therefore froze the model with the original ($0.004$, threshold one) coefficient for the entire manuscript --- this is the basin in which the therapeutic surface of Section~\ref{sec:r-q4r} is mapped, and which mirrors the classical amyloid-driven phenotype of the disease.

The existence of a redox-collapse basin orthogonal to the classical amyloid basin is a separate prediction that we leave to future work. Biologically, it corresponds to the literature on chronic glutathione depletion as a driver of neurodegeneration via mechanisms partly orthogonal to amyloid-$\beta$~\cite{cassano2020}. Methodologically, it underlines the principle that even surgical one-line edits to a multi-stable biological Petri net can shift basin boundaries, and that re-baselining must follow every patch.

\subsection{Limitations}

Several caveats apply. We adopt the convention that one marking unit is equivalent to one micromolar for the inflammatory dose-response interpretation; direct translation to clinical pharmacology requires explicit pharmacokinetic--pharmacodynamic calibration against measured brain and plasma concentrations. The model does not include tau hyperphosphorylation or neurofibrillary tangles, despite cannabidiol's documented effects on the tau axis~\cite{esposito2006tau}. The brain is treated as a single well-mixed compartment, omitting region-specific differences in cannabidiol distribution, microglial density, or vulnerability. The age-dependent mechanism switch reported in earlier work~\cite{simao2026shypn} was a finding of the previous model and has not been re-validated against the frozen v3 baseline; the age compartment remains in the topology and a future age-axis sweep would close this gap. The model omits cytochrome P450 3A4 cannabidiol metabolism and therefore does not predict the inverted-U dose response sometimes observed in high-dose studies. Finally, amyloid precursor protein processing is treated kinetically without explicit endolysosomal pH compartmentalisation.

Two additional caveats are particular to the present manuscript. The seven-day integration horizon depletes the neuron-health pool below the resolution of stochastic noise across all parameter combinations explored, regardless of treatment, so neuron-health endpoints are not separable across the factorial at this horizon. We therefore omit neuron-health endpoints from the present manuscript and report biomarker-level readouts (NF$\kappa$B p65, amyloid-$\beta$ oligomer, reactive oxygen species, reduced glutathione) as the principal observables; a horizon-shortening study, or a re-parameterisation of the neurotoxicity-versus-neuroprotection balance to operate on a longer horizon without floor saturation, is left for future work and is not expected to alter the qualitative findings reported here. Separately, the amyloid-$\beta$ production rate currently embeds the disease-severity parameter directly into its rate expression --- a Pattern-A leak relative to the canonical experiment-plan bridge of parameter $\rightarrow$ event $\rightarrow$ spatial signal $\rightarrow$ rate. The leak takes the form of a $(1 + 2 \cdot \text{severity})$ multiplier and is identical to its default at severity zero, so the qualitative findings reported here are unaffected; a topological refactor that promotes the multiplier to a spatial-signal place written by an event at $t = 0$ is scheduled for a v4 model and will be reported separately.

\subsection{Reproducibility and Provenance}

All numerical results are reproducible from the model file at git tag \texttt{v3-manuscript-frozen} (sha256 prefix \texttt{a32bf9a8b2dd2afc}) and engine commit \texttt{7fd06fd8} via the three sweep configurations \texttt{sweep\_config.Q4r-final.json}, \texttt{sweep\_config.Q5-final.json}, and \texttt{sweep\_config.Q3-final.json}, committed alongside this manuscript. Each results directory carries a \texttt{provenance.json} (client and server commit, branch, dirty-tree status, dispatch timestamp) and a byte-identical \texttt{model\_snapshot.shy}; per-condition CSV summaries and per-replicate trajectory archives are stored under \texttt{workspace/projects/canabidiol/experiments/results/}.

% ============================================================================
% CONCLUSIONS
% ============================================================================
\section{Conclusions}

We present an in silico laboratory for cannabidiol neuroprotection in Alzheimer's disease, built on a hybrid stochastic--continuous biological Petri net (forty-three places, forty-eight transitions, one hundred and four arcs, forty-three discrete events) integrated over a seven-day horizon under realistic twenty-seven-event maintenance dosing. Three findings emerge from sweeps run against a single frozen baseline. First, the untreated disease cascade is strict-monotone in severity across every measured biomarker (NF$\kappa$B p65 from $0.65$ to $2.35$, amyloid-$\beta$ oligomer from $8.8$ to $117$, reduced glutathione depleted from $80$ to $5$), with conservation invariants preserved exactly --- confirming that the experiment-plan layer perturbs only the biological initial condition and that the subsequent therapeutic findings are emergent properties of the topology. Second, cannabidiol suppresses NF$\kappa$B p65 along a Hill curve with half-maximal inhibitory concentration of $0.08$ marking units and Hill coefficient near $1.65$, essentially independent of disease severity, and clears amyloid-$\beta$ oligomer to zero over the seven-day horizon at every non-zero maintenance dose; the operative dose lies an order of magnitude below the one-to-thirty micromolar range of historical preclinical studies. Third, the system carries an intrinsic glutathione-mediated bistable cusp on the oxidative axis whose location is itself a function of treatment: the bifurcation lies at moderate-to-severe disease in the untreated cascade and migrates to mild disease under maintenance cannabidiol, with the bimodal redox stratum shifting by approximately one severity step. Cannabidiol therefore does not abolish basin selection on the redox axis but reassigns which severity stratum sits on the cusp, motivating a glutathione-restoring co-therapy --- N-acetylcysteine the natural candidate --- and stratification of clinical cohorts by oxidative-stress biomarkers in addition to amyloid load. These predictions are accompanied by explicit refutation criteria and supporting protocols; the full pipeline (frozen model, three sweep configurations, per-run provenance, and per-replicate trajectories) is archived for independent reproduction.

% ============================================================================
% MATTER
% ============================================================================
\section*{Acknowledgments}

The author thanks the SHyPN development community and UFSC for computational resources.

\section*{Funding}

\todo{Funding sources}

\section*{Data Availability}

All model files, sweep configurations, run provenance, and per-condition CSV data are archived in the public repository at \todo{repo URL} under git tag \texttt{v3-manuscript-frozen}. Each results directory contains \texttt{provenance.json} and a byte-identical \texttt{model\_snapshot.shy} for the exact state simulated.

\section*{Author Contributions}

E.S.\ conceived the study, developed the SHyPN framework, designed and dispatched the sweeps, analysed the results, and wrote the manuscript.

\section*{Competing Interests}

The author declares no competing interests.

\bibliographystyle{plos2015}
\bibliography{references}

\end{document}
